% =============================================================================
%  Etapa I - Propuesta y Definicion del Problema
%  Proyecto Semestral de Data Science
%
%  Compilar con:  pdflatex etapa1.tex   (dos veces, para el indice)
%
%  Paquetes usados: geometry, babel, booktabs, longtable, array, amsmath,
%  hyperref, enumitem, xcolor, titlesec, fancyhdr.
%  Todos vienen en texlive-latex-recommended y texlive-latex-extra, y estan
%  disponibles en Overleaf sin instalar nada.
% =============================================================================

\documentclass[11pt,a4paper]{article}

\usepackage[T1]{fontenc}
\usepackage[utf8]{inputenc}
\usepackage[spanish]{babel}

% --- Margenes -----------------------------------------------------------------
\usepackage[top=2.5cm,bottom=2.5cm,left=3cm,right=3cm]{geometry}

% --- Tablas -------------------------------------------------------------------
\usepackage{booktabs}   % \toprule \midrule \bottomrule
\usepackage{longtable}  % tablas que cruzan de pagina
\usepackage{array}      % columnas p{} con formato

% --- Matematicas --------------------------------------------------------------
\usepackage{amsmath}

% --- Listas mas compactas -----------------------------------------------------
\usepackage{enumitem}
\setlist{itemsep=2pt,parsep=0pt,topsep=4pt}

% --- Color y titulos ----------------------------------------------------------
\usepackage{xcolor}
\definecolor{azul}{HTML}{1F4E79}

\usepackage{titlesec}
\titleformat{\section}{\Large\bfseries\color{azul}}{\thesection.}{0.6em}{}
\titleformat{\subsection}{\large\bfseries\color{azul}}{\thesubsection}{0.6em}{}
\titlespacing*{\section}{0pt}{18pt}{8pt}
\titlespacing*{\subsection}{0pt}{12pt}{6pt}

% --- Numeracion de paginas ----------------------------------------------------
\usepackage{fancyhdr}
\pagestyle{fancy}
\fancyhf{}
\fancyfoot[C]{\thepage}
\renewcommand{\headrulewidth}{0pt}

% --- Enlaces ------------------------------------------------------------------
\usepackage[hidelinks]{hyperref}

\begin{document}

% =============================================================================
%  PORTADA
% =============================================================================
\begin{titlepage}
\centering
\vspace*{3cm}

{\Huge\bfseries\color{azul} Riesgo de siniestros de tránsito\\[0.3em]
en carreteras de Chile\par}

\vspace{1.2cm}
{\Large\bfseries Etapa I --- Propuesta y Definición del Problema\par}

\vspace{2.5cm}
{\large Proyecto Semestral de Data Science\par}

\vspace{0.6cm}
{\large 11 de septiembre de 2026\par}

\vspace{3cm}
\begin{tabular}{l}
\textbf{Integrantes del equipo:}\\[0.4em]
% -------- COMPLETAR AQUÍ (4 a 5 integrantes) --------
Nombre Apellido\\
Nombre Apellido\\
Nombre Apellido\\
Nombre Apellido\\
% ---------------------------------------------------
\end{tabular}

\vfill
{\small
\textbf{Repositorio:} \url{https://github.com/martinnmg1809/siniestros-transito-chile}\\[0.3em]
\textbf{Prototipo funcional:} \url{https://martinnmg1809.github.io/siniestros-transito-chile}
\par}

\end{titlepage}

% =============================================================================
%  ÍNDICE
% =============================================================================
\tableofcontents
\newpage

% =============================================================================
\section{Resumen de la propuesta}

En Chile mueren cerca de 1.400 personas al año en siniestros de tránsito. La
respuesta institucional se apoya en mapas de puntos críticos que se publican una
vez al año y que, por construcción, describen \textbf{dónde} ocurrieron los
siniestros pero no \textbf{cuándo} el riesgo se eleva.

Este proyecto propone estimar el \textbf{número esperado de siniestros por tramo
de carretera y por hora}, combinando el registro histórico georreferenciado de
Carabineros ---publicado por CONASET--- con variables de contexto conocidas por
anticipado: pronóstico de precipitación, hora del día, día de la semana y
calendario de feriados.

El problema se formula como una \textbf{regresión de conteos sobre datos
espacio-temporales}, no como una clasificación. Esa decisión es el eje
conceptual de la propuesta y se justifica en la sección~\ref{sec:enfoque}: un
siniestro en un tramo y hora concretos es un evento con probabilidad cercana al
0,13\,\%, de modo que cualquier clasificador binario alcanzaría una exactitud
superior al 99\,\% prediciendo siempre «no habrá siniestro», sin aportar
información alguna.

El entregable final es doble: un modelo validado temporalmente y una interfaz
que traduce sus estimaciones a una decisión operativa concreta: \emph{qué tramos
vigilar en las próximas doce horas}.

% =============================================================================
\section{Contexto del problema}

\subsection{Dominio de aplicación}

El proyecto se sitúa en la \textbf{seguridad vial interurbana chilena},
específicamente en la red de carreteras y autopistas concesionadas administradas
por el Ministerio de Obras Públicas y fiscalizadas por Carabineros de Chile.

La magnitud del fenómeno es la siguiente, según los informes nacionales de
CONASET:

\begin{center}
\begin{tabular}{lrrr}
\toprule
\textbf{Año} & \textbf{Siniestros} & \textbf{Fallecidos} & \textbf{Lesionados}\\
\midrule
2023 & 78.238 & 1.635 & 45.679\\
2024 & 75.653 & 1.439 & 42.400\\
\bottomrule
\end{tabular}
\end{center}

La cifra de fallecidos bajó un 12\,\% entre 2023 y 2024, pero se mantiene en un
orden de magnitud que sitúa al tránsito entre las principales causas de muerte
por causas externas en el país.

\subsection{Por qué la red interurbana}

De los 436.521 siniestros georreferenciados que CONASET publica para el período
2020--2025, 96.290 ocurren en \textbf{zona rural} y 340.231 en zona urbana.
Aunque la zona urbana concentra más eventos, la red interurbana presenta tres
características que la hacen un dominio más tratable y más relevante para este
proyecto:

\begin{enumerate}
  \item \textbf{La letalidad es mayor.} Las velocidades de operación en
  carretera son sustancialmente más altas, de modo que un mismo número de
  siniestros produce más fallecidos.

  \item \textbf{La geometría es unidimensional.} Una carretera se describe por
  su kilometraje, lo que permite una referenciación lineal precisa y una unidad
  de análisis natural: el tramo. En zona urbana el evento ocurre en una red de
  intersecciones cuya topología es mucho más compleja de modelar.

  \item \textbf{La intervención es concentrable.} Un patrullaje o una
  señalización variable desplegada en un tramo específico de carretera cubre la
  totalidad del flujo que circula por él, algo que no ocurre en una malla urbana
  con múltiples rutas alternativas.
\end{enumerate}

\subsection{La problemática concreta}

CONASET publica hoy capas de \textbf{puntos críticos} con periodicidad anual.
Esta información es valiosa pero tiene una limitación estructural: es
\textbf{estática}. Un tramo aparece marcado como crítico porque históricamente
acumuló siniestros, y esa marca es la misma un martes a las 3 de la madrugada
con la calzada seca que un viernes a las 18:00 bajo lluvia.

La consecuencia práctica es que la asignación de recursos de fiscalización no
puede priorizarse en el tiempo. Se sabe \emph{dónde} mirar, pero no
\emph{cuándo} mirar con más atención.

Esta brecha existe pese a que las variables que modulan el riesgo son
\textbf{conocidas con antelación}: el pronóstico meteorológico está disponible
con horas o días de anticipación, y la hora, el día de la semana y el calendario
de feriados son deterministas.

% =============================================================================
\section{Justificación}

\subsection{Relevancia}

La justificación del proyecto no reside en predecir siniestros individuales
---algo que la sección~\ref{sec:problema} descarta explícitamente--- sino en
\textbf{mejorar la eficiencia con que se asignan recursos escasos de
fiscalización}.

El argumento es cuantificable. Si un cuerpo de fiscalización puede cubrir una
fracción limitada de la red en un turno determinado, la pregunta operativa es
qué fracción de los siniestros reales quedaría cubierta por esa asignación.
Nuestra evidencia preliminar (sección~\ref{sec:evidencia}) indica que:

\begin{itemize}
  \item Con la información estática actual ---solo la ubicación histórica---
  vigilar el 5\,\% de las horas-tramo disponibles cubre el \textbf{14,9\,\%} de
  los siniestros reales del año siguiente.
  \item Incorporando hora, día, feriado y pronóstico de lluvia, esa cobertura
  sube al \textbf{19,2\,\%}.
\end{itemize}

La diferencia es una mejora relativa del 28\,\% en la eficiencia de cobertura,
obtenida \textbf{sin desplegar un solo recurso adicional}: únicamente
reordenando cuándo se despliegan los que ya existen.

\subsection{Beneficiarios e impacto esperado}

\begin{center}
\begin{tabular}{p{0.32\textwidth}p{0.58\textwidth}}
\toprule
\textbf{Actor} & \textbf{Uso concreto de la herramienta}\\
\midrule
Carabineros (SIAT y unidades de tránsito) & Priorización temporal de patrullajes
y controles preventivos\\
\addlinespace
Concesionarias viales & Anticipación del despliegue de asistencia en ruta y
señalización variable\\
\addlinespace
CONASET / MOP & Complemento dinámico a los mapas anuales de puntos críticos\\
\addlinespace
Conductores & Información de riesgo por tramo antes de emprender un viaje\\
\bottomrule
\end{tabular}
\end{center}

\subsection{Por qué este proyecto es viable ahora}

Tres condiciones habilitantes se cumplen simultáneamente, y ninguna de ellas era
cierta hace algunos años:

\begin{enumerate}
  \item \textbf{Los datos están abiertos y georreferenciados.} CONASET publica
  los registros de Carabineros en servicios ArcGIS de acceso público, con
  coordenadas y kilometraje de ruta, sin necesidad de convenios ni credenciales.

  \item \textbf{El clima histórico y pronosticado es accesible sin costo.}
  Existen APIs de reanálisis meteorológico con resolución horaria y cobertura
  global que permiten reconstruir las condiciones de cada siniestro pasado y
  consultar el pronóstico de las próximas horas.

  \item \textbf{La entrega no requiere infraestructura.} El modelo, una vez
  ajustado, se reduce a un conjunto reducido de coeficientes, lo que permite
  publicar el resultado como una página estática sin servidor ni base de datos.
\end{enumerate}

% =============================================================================
\section{Definición del problema}
\label{sec:problema}

\subsection{Pregunta central}

\begin{quote}
\textbf{¿Cuántos siniestros de tránsito se espera que ocurran en cada tramo de
una carretera chilena durante las próximas doce horas, y qué tramos concentran
el riesgo, dadas la hora, el día de la semana, el calendario de feriados y el
pronóstico de precipitación?}
\end{quote}

\subsection{Unidad de análisis}

La unidad de análisis es la \textbf{celda tramo $\times$ hora}: un segmento de
carretera de longitud fija (5\,km como caso base) durante una hora calendario
determinada. Para la Ruta 5 Sur, el período 2020--2024 genera 251 tramos
$\times$ 43.848 horas, es decir, aproximadamente \textbf{11,0 millones de
celdas}.

\subsection{Qué se busca resolver, y qué no}

Delimitar esto es esencial para la coherencia del proyecto, y constituye la
principal decisión conceptual de la propuesta.

\medskip
\noindent\textbf{El proyecto SÍ busca responder:}
\begin{itemize}
  \item Cuántos siniestros esperar en una ruta completa en una ventana de doce
  horas.
  \item Qué tramos concentran el riesgo dentro de esa ventana, y en qué
  proporción respecto del tramo promedio.
  \item Cuánto eleva el riesgo cada factor de contexto: lluvia, hora punta, día
  de salida a un fin de semana largo.
\end{itemize}

\medskip
\noindent\textbf{El proyecto NO busca responder:}
\begin{itemize}
  \item Si un tramo y hora específicos tendrán un siniestro. La tasa base es de
  \textbf{0,133\,\% por tramo-hora}; en toda la Ruta 5 Sur se esperan
  \textbf{4,0 siniestros en doce horas} repartidos a lo largo de 1.226\,km. Una
  afirmación del tipo «mañana a las 15:20 en el km 193 habrá una colisión» no es
  alcanzable con estos datos, ni con ningún otro conjunto de datos
  observacionales de esta naturaleza.
  \item Las causas del siniestro. El registro contiene una causa asignada por el
  funcionario que concurrió al lugar, pero el diseño observacional no permite
  inferencia causal.
\end{itemize}

\medskip
Esta delimitación no es una limitación del método, sino una propiedad del
fenómeno: los siniestros son eventos raros en el espacio-tiempo. Reconocerlo
explícitamente es lo que permite formular un objetivo que sí es alcanzable y
verificable.

% =============================================================================
\section{Enfoque propuesto}
\label{sec:enfoque}

\subsection{Tipo de problema de Ciencia de Datos}

El problema se formula como una \textbf{regresión de conteos sobre datos
espacio-temporales}, mediante un modelo lineal generalizado de Poisson con
enlace logarítmico y término de exposición (\emph{offset}).

Formalmente, para la celda correspondiente al tramo $s$ y la hora $t$:

\begin{equation*}
Y_{st} \sim \mathrm{Poisson}(\lambda_{st})
\end{equation*}
\begin{equation*}
\log \lambda_{st} = \alpha_s + \beta_{h(t)} + \gamma_{d(t)} + \delta_{m(t)}
+ \phi_{f(t)} + \psi_{c(s,t)} + \log E_{st}
\end{equation*}

donde $\alpha_s$ es el efecto del tramo, $\beta$ la hora del día, $\gamma$ el
día de la semana, $\delta$ el mes, $\phi$ la posición respecto de un feriado,
$\psi$ la categoría de precipitación y $E_{st}$ la exposición en horas.

\subsection{Justificación del enfoque, frente a las alternativas}

\paragraph{Por qué no clasificación binaria.} Sería el enfoque intuitivo
---«¿habrá o no habrá siniestro?»--- y es el que descartamos con mayor
convicción. Con una tasa base de 0,133\,\%, un clasificador que responda siempre
«no» alcanza un 99,87\,\% de exactitud y un AUC que puede parecer razonable, sin
contener información útil. Además, la clasificación descarta información: dos
celdas con riesgo esperado de 0,001 y 0,05 siniestros son ambas «clase
negativa», cuando su riesgo difiere en un factor de cincuenta. La magnitud es
precisamente lo que interesa a quien asigna recursos.

\paragraph{Por qué no series de tiempo puras.} Un modelo ARIMA o similar sobre
el conteo agregado de la ruta capturaría la estacionalidad temporal pero
\textbf{perdería la dimensión espacial}, que es donde nuestra evidencia
preliminar indica que reside la mayor parte de la señal. La pregunta operativa
no es cuántos siniestros habrá en la ruta, sino en qué tramo.

\paragraph{Por qué no clustering.} Es un enfoque no supervisado y aquí sí existe
una variable respuesta observada: el conteo de siniestros. Puede ser útil como
análisis exploratorio complementario para agrupar tramos con perfiles de riesgo
similares, pero no responde la pregunta central.

\paragraph{Por qué Poisson y no otra regresión.} La variable respuesta es un
conteo entero no negativo, con gran mayoría de ceros, y el término de exposición
permite modelar celdas con distinta cantidad de horas observadas. Se contrastará
contra una \textbf{binomial negativa} para evaluar sobredispersión, y contra un
\textbf{\emph{gradient boosting} con objetivo de Poisson} para capturar
interacciones no aditivas.

\subsection{Marco metodológico}

Se adopta \textbf{CRISP-DM}, por tres razones concretas:

\begin{enumerate}
  \item Su primera fase es \emph{Business Understanding}, que en este proyecto
  es la fase crítica: la decisión de formular el problema como conteos y no como
  clasificación es una decisión de entendimiento del dominio, no técnica.

  \item Es \textbf{iterativo entre preparación de datos y modelado}, lo cual es
  necesario aquí: los problemas de calidad del kilometraje descritos en la
  sección~\ref{sec:limitaciones} solo se hicieron visibles al intentar modelar.

  \item Su fase de \emph{Deployment} es un entregable explícito, coherente con
  que el resultado del proyecto sea una interfaz operativa y no solo un modelo.
\end{enumerate}

\subsection{Estrategia de validación}

Dos decisiones metodológicas definen la validación:

\paragraph{Partición temporal, no aleatoria.} Se entrena con 2020--2023 y se
evalúa con 2024. Una partición aleatoria sobre datos con autocorrelación
espacio-temporal filtraría información del futuro al conjunto de entrenamiento e
inflaría artificialmente las métricas.

\paragraph{Métrica operativa por sobre métrica estadística.} Además de la
desvianza de Poisson, la métrica principal es la \textbf{curva de captura}:
\emph{si se dedica el N\,\% de las horas-tramo disponibles a vigilar las de
mayor riesgo estimado, qué porcentaje de los siniestros reales del período de
prueba queda cubierto}. Esta métrica traduce directamente el desempeño del
modelo a la decisión que se busca apoyar.

\paragraph{El baseline exigente es el mapa estático.} No basta con superar al
azar. El modelo debe superar a un modelo que use \textbf{únicamente la ubicación
del tramo}, porque ese es el estado del arte actual: es lo que ya publica
CONASET. Si añadir hora, clima y feriados no mejora la captura fuera de muestra,
esa capa está añadiendo ruido y no debe publicarse.

% =============================================================================
\section{Objetivos del proyecto}

\subsection{Objetivo general}

Desarrollar y validar un sistema reproducible que estime el número esperado de
siniestros de tránsito por tramo y hora en carreteras de Chile, en función de
variables de contexto conocidas por anticipado, y comunicar sus estimaciones en
una interfaz que apoye la priorización temporal y espacial de recursos de
fiscalización.

\subsection{Objetivos específicos}

Cada objetivo se acompaña de una \textbf{meta verificable} al cierre del
proyecto.

\begin{longtable}[c]{p{0.46\textwidth}p{0.46\textwidth}}
\toprule
\textbf{Objetivo específico} & \textbf{Meta verificable}\\
\midrule
\endfirsthead
\toprule
\textbf{Objetivo específico} & \textbf{Meta verificable}\\
\midrule
\endhead
\bottomrule
\endlastfoot

\textbf{OE1.} Consolidar un conjunto de datos reproducible de siniestros con
ruta y kilometraje validados, a partir de las fuentes abiertas de CONASET.
& Al menos 80.000 siniestros con ruta asignada; $\geq$ 95\,\% con coordenadas
válidas; pipeline ejecutable de extremo a extremo con un solo comando.\\
\addlinespace

\textbf{OE2.} Diagnosticar y corregir los problemas de calidad del campo de
kilometraje, documentando el criterio aplicado.
& Correlación $\geq$ 0,95 entre el kilometraje declarado y la posición
geográfica proyectada, en toda ruta que se incorpore al modelo.\\
\addlinespace

\textbf{OE3.} Cuantificar el efecto de las condiciones meteorológicas, la hora,
el día de la semana y los feriados sobre la tasa de siniestros.
& Efectos multiplicativos estimados con intervalo de confianza para cada
variable, sobre la Ruta 5 Sur y al menos cuatro rutas adicionales.\\
\addlinespace

\textbf{OE4.} Ajustar un modelo de conteos que supere de forma demostrable al
baseline puramente espacial en validación temporal.
& Captura en el 5\,\% superior de horas-tramo al menos 3 puntos porcentuales por
encima del modelo de solo ubicación, evaluado sobre 2024.\\
\addlinespace

\textbf{OE5.} Comparar al menos tres especificaciones de modelo y seleccionar
una con criterio explícito.
& Tabla comparativa Poisson / binomial negativa / \emph{gradient boosting} con
desvianza y curva de captura fuera de muestra.\\
\addlinespace

\textbf{OE6.} Extender el análisis a múltiples rutas, aplicando el criterio de
validación como filtro de publicación.
& Al menos 5 rutas publicadas que superen el baseline estático, y documentación
explícita de las rutas descartadas con su motivo.\\
\addlinespace

\textbf{OE7.} Entregar una interfaz operativa que muestre el riesgo de las
próximas 12 horas por tramo, en dos vistas complementarias.
& Interfaz accesible públicamente, que se actualice con el pronóstico vigente y
que degrade de forma controlada si la fuente meteorológica no responde.\\
\addlinespace

\textbf{OE8.} Documentar de forma explícita las limitaciones del sistema y la
incertidumbre de sus estimaciones.
& Sección de limitaciones que aborde al menos: ausencia de flujo vehicular
horario, subregistro del registro policial y error de medición de la variable
meteorológica.\\

\end{longtable}

% =============================================================================
\section{Evidencia preliminar de factibilidad}
\label{sec:evidencia}

Se realizó una exploración inicial para verificar que los datos sostienen la
propuesta antes de comprometerse con ella. Los resultados de esta sección son
\textbf{preliminares} y se presentan como sustento de viabilidad, no como
resultados finales del proyecto.

\subsection{Disponibilidad de datos}

\begin{center}
\begin{tabular}{p{0.30\textwidth}p{0.60\textwidth}}
\toprule
\textbf{Fuente} & \textbf{Contenido verificado}\\
\midrule
CONASET --- ArcGIS Hub & 436.521 siniestros georreferenciados 2020--2025, con
fecha, hora, tipo, causa, ruta y kilometraje\\
\addlinespace
CONASET --- capa de feriados & Calendario 2019--2024, que distingue el día de
salida, los días intermedios y el de regreso de cada fin de semana largo\\
\addlinespace
CONASET --- capas regionales 2023 & Estado atmosférico y condición de calzada
observados en terreno\\
\addlinespace
Open-Meteo & Precipitación horaria, histórica y pronosticada, sin credenciales\\
\addlinespace
MOP --- Plan Nacional de Censos & 863 puntos censales con Tránsito Medio Diario
Anual por rol de ruta\\
\bottomrule
\end{tabular}
\end{center}

\medskip
De los 82.026 siniestros con ruta asignada, la \textbf{Ruta 5 Sur} concentra
14.282 tras limpieza, 2,5 veces más que la siguiente ruta con más registros.

\subsection{Resultados preliminares del modelo}

Ajustando el modelo descrito en la sección~\ref{sec:enfoque} sobre la Ruta 5
Sur, con entrenamiento en 2020--2023 y prueba sobre los 2.585 siniestros de
2024:

\begin{center}
\begin{tabular}{lrr}
\toprule
\textbf{Modelo} & \textbf{Captura top 5\,\%} & \textbf{Captura top 10\,\%}\\
\midrule
Tasa constante (azar)      & 2,1\,\%  & 5,7\,\%\\
Solo ubicación del tramo   & 14,9\,\% & 26,3\,\%\\
\textbf{Modelo completo}   & \textbf{19,2\,\%} & \textbf{30,6\,\%}\\
\bottomrule
\end{tabular}
\end{center}

\medskip
\noindent\textbf{Efectos multiplicativos estimados} sobre la celda promedio:

\begin{center}
\begin{tabular}{lr}
\toprule
\textbf{Variable} & \textbf{Efecto}\\
\midrule
Lluvia superior a 0,5\,mm/h              & \textbf{2,13$\times$}\\
Llovizna (0,1--0,5\,mm/h)                & 1,31$\times$\\
Hora punta (18:00)                       & 1,68$\times$\\
Día de salida a un fin de semana largo   & 1,19$\times$\\
Día feriado intermedio                   & 0,93$\times$\\
Madrugada (02:00)                        & 0,52$\times$\\
Día de la semana (rango completo)        & 0,93$\times$ -- 1,10$\times$\\
\bottomrule
\end{tabular}
\end{center}

\medskip
Tres hallazgos preliminares orientan el trabajo restante:

\begin{enumerate}
  \item \textbf{La componente espacial aporta la mayor parte de la señal.} El
  mapa estático ya alcanza el 14,9\,\% de captura. La capa dinámica añade 4,3
  puntos. Es una mejora real y verificable, pero conviene declararla en su
  magnitud correcta y no sobrevenderla.

  \item \textbf{La precipitación es el factor de contexto más fuerte}, por
  encima de la hora punta.

  \item \textbf{El día de salida a un feriado es más riesgoso que el feriado
  mismo} (1,19$\times$ frente a 0,93$\times$), lo que sugiere que la variable
  relevante es el desplazamiento y no la fecha feriada en sí.
\end{enumerate}

\subsection{Una validación que arrojó resultado negativo}

Se evaluó si el reanálisis meteorológico reproduce la observación de estado
atmosférico registrada por Carabineros, contrastando 6.096 siniestros de la
Región del Biobío en 2023. El acuerdo resultó \textbf{débil}: índice kappa de
Cohen de 0,262 y sensibilidad de 47,2\,\%, es decir, el reanálisis no detecta
más de la mitad de las lluvias registradas en terreno. La causa es estructural:
el reanálisis promedia sobre celdas de aproximadamente 25\,km, y la
precipitación en la zona centro-sur del país es localizada.

Este resultado obliga a una distinción conceptual que se incorpora al proyecto:
\textbf{reproducir la etiqueta observada y predecir el conteo de siniestros son
dos preguntas distintas}. El reanálisis falla en la primera y, sin embargo, la
precipitación resulta el predictor más fuerte del modelo. El error de medición
atenúa el coeficiente estimado hacia 1, de modo que el efecto real de la lluvia
es probablemente mayor que el 2,13$\times$ estimado.

% =============================================================================
\section{Datos, riesgos y limitaciones}
\label{sec:limitaciones}

\subsection{Problemas de calidad identificados}

\paragraph{Unidades mezcladas en el kilometraje.} En varias rutas una fracción
de los registros expresa la ubicación en metros: el valor 84.300 de la Ruta 68
corresponde en realidad al km 84,3. Sin embargo, la Ruta 5 Sur alcanza
legítimamente el km 1.226, por lo que una corrección global destruiría sus
datos. El criterio adoptado es \textbf{por ruta}: se elige la transformación que
maximiza la correlación entre el kilometraje declarado y la posición geográfica
proyectada sobre el eje de la ruta. En la Ruta 68 esa correlación pasa de 0,139 a
0,999; en la Ruta 5 Sur se mantiene en 0,998 sin transformar.

\paragraph{Campo de ruta incompleto.} Los registros de 2025 existen en la base
pero no tienen el campo de ruta poblado, lo que limita el período utilizable a
2020--2024.

\subsection{Limitaciones que se declararán en el resultado final}

\begin{center}
\begin{tabular}{p{0.38\textwidth}p{0.52\textwidth}}
\toprule
\textbf{Limitación} & \textbf{Consecuencia}\\
\midrule
\textbf{Ausencia de flujo vehicular horario.} El censo del MOP entrega un
promedio anual por punto censal.
& El modelo no puede separar «hay más vehículos a esta hora» de «conducir a esta
hora es más peligroso por vehículo». Para priorizar recursos esto es
indiferente; para una afirmación causal sobre el riesgo, no.\\
\addlinespace
\textbf{Subregistro del registro policial.} La fuente son partes de Carabineros.
& Los siniestros sin lesionados y aquellos donde no concurrió personal policial
están subrepresentados.\\
\addlinespace
\textbf{Error de medición de la variable meteorológica.}
& Los coeficientes de precipitación están atenuados; el efecto real es
probablemente mayor.\\
\addlinespace
\textbf{Modelos no comparables entre rutas.} Cada ruta tiene su propio modelo y
su propio nivel basal.
& Un factor de «2,0$\times$» en una ruta no equivale al mismo factor en otra.\\
\bottomrule
\end{tabular}
\end{center}

\subsection{Riesgos del proyecto y mitigación}

\begin{center}
\begin{tabular}{p{0.42\textwidth}p{0.48\textwidth}}
\toprule
\textbf{Riesgo} & \textbf{Mitigación}\\
\midrule
Rutas con pocos datos producen modelos que son puro ruido
& El criterio de publicación exige superar al baseline estático en validación
temporal; las rutas que no lo cumplen se documentan pero no se publican\\
\addlinespace
Dependencia de una API externa para el pronóstico
& La interfaz calcula con calzada seca y lo advierte cuando la fuente no
responde; el modelo sin clima conserva la mayor parte de su capacidad\\
\addlinespace
Sobreajuste por el número de parámetros de tramo
& Regularización L2 sobre los efectos de tramo y validación estrictamente
temporal\\
\bottomrule
\end{tabular}
\end{center}

% =============================================================================
\section{Plan de trabajo}

El plan se alinea con las cuatro entregas de la asignatura y con las fases de
CRISP-DM.

\begin{center}
\begin{tabular}{p{0.22\textwidth}p{0.24\textwidth}p{0.42\textwidth}}
\toprule
\textbf{Etapa} & \textbf{Fase CRISP-DM} & \textbf{Actividades principales}\\
\midrule
\textbf{I} --- Propuesta (11 sep)
& \emph{Business Understanding}
& Definición del problema, formulación como regresión de conteos, objetivos
medibles\\
\addlinespace
\textbf{II} --- Avance (22--26 oct)
& \emph{Data Understanding} y \emph{Data Preparation}
& Consolidación del dataset, corrección del kilometraje, análisis exploratorio,
primer modelo Poisson con resultados preliminares\\
\addlinespace
\textbf{III} --- Presentación técnica (16--19 nov)
& \emph{Modeling} y \emph{Evaluation}
& Comparación Poisson / binomial negativa / \emph{gradient boosting}, selección
con criterio explícito, validación temporal, extensión a múltiples rutas\\
\addlinespace
\textbf{IV} --- Feria de Proyectos (2 dic)
& \emph{Deployment}
& Interfaz operativa, póster, video resumen y comunicación a audiencia no
especializada\\
\bottomrule
\end{tabular}
\end{center}

% =============================================================================
\section{Fuentes}

\begin{itemize}
  \item CONASET. \emph{Portal de mapas y datos abiertos.}
  \url{https://mapas-conaset.opendata.arcgis.com}
  \item CONASET. \emph{Informe Nacional de Siniestros de Tránsito en Chile
  2024.}
  \item CONASET. \emph{Informe Nacional de Siniestros de Tránsito en Chile
  2023.}
  \item Carabineros de Chile. \emph{Anuario Estadístico de Accidentes en el
  Tránsito} (sistema SIEC 2).
  \item Ministerio de Obras Públicas, Dirección de Vialidad. \emph{Plan Nacional
  de Censos Viales.}
  \item Open-Meteo. \emph{Historical Weather API} y \emph{Forecast API}
  (reanálisis ERA5).
  \item Dirección Meteorológica de Chile. \emph{Portal de Servicios Climáticos.}
\end{itemize}

\end{document}
